\section{Conclusions}
\label{sec:conclusions}

This paper estimates the contribution of informal work to Colombian output from 1993 to 2026. In a two-sector DSGE model, the informal output share depends on three quantities: the proportion of workers employed informally, their value added per worker relative to formal workers, and the capital share of formal production. The first is observed in the informality rate, the second is measured with a new quarterly series of relative earnings built from household-survey microdata, and the third is estimated from Colombian data.

The identification analysis shows why the relative-earnings series is needed. Aggregate national accounts and labor-market series do not determine the size of the informal sector on their own: likelihood profiles, prior--posterior comparisons and simulations with known parameters all point in that direction. With relative earnings, the quarterly path of the informal share is identified, invariant to the model's steady state and recovered precisely in simulations. The distinction between the path, which the data identify, and the steady-state parameter, which they do not need to identify, is the key to estimating the informal economy with a DSGE model.

Under the baseline specification, the informal share falls from about 30\% of GDP in 1993 to 20\% in 2026Q2, while labor informality falls from about 69\% to 54\%. Informal value added per worker is about one-fifth of formal value added per worker, which explains how a majority of workers can be informal while producing a fifth of output. The estimates lie within the range of published Colombian series and close to the most recent MIMIC estimates.

The model delivers more than the informal share: quarterly informal and formal GDP, relative productivity, the informality rate net of measurement error and the structural shocks behind them, all of which are published with the paper. The data and code allow the series to be updated every quarter and extended as new information becomes available, in particular direct measures of the value added of informal units and observed policy instruments that would turn the model into a tool for evaluating formalization policies.
